% BHU PYQ -- Professional Exam Template
\documentclass[11pt,a4paper]{article}

\usepackage[a4paper,top=2.5cm,bottom=2.5cm,left=2.8cm,right=2.8cm,headheight=14pt]{geometry}
\usepackage{amsmath,amssymb}
\usepackage[T1]{fontenc}
\usepackage[utf8]{inputenc}
\usepackage{lmodern}
\usepackage{microtype}
\usepackage{fancyhdr}
\usepackage{enumitem}
\usepackage{listings}
\usepackage{hyperref}
\usepackage{lastpage}
\usepackage{parskip}

\hypersetup{colorlinks=true,urlcolor=black,linkcolor=black,
  pdftitle={BVCA-402: Probability and Statistics -- BHU PYQ}}

\pagestyle{fancy}
\fancyhf{}
\renewcommand{\headrulewidth}{0.4pt}
\renewcommand{\footrulewidth}{0.4pt}
\fancyhead[L]{\small   \textit{Banaras Hindu University}}
\fancyhead[R]{\small   \textit{BVCA-402: Probability and Statistics}}
\fancyfoot[C]{\small Page  \thepage\ of \pageref{LastPage}}

\lstset{
  basicstyle=\small tfamily,
  frame=single,
  breaklines=true,
  columns=flexible,
  keepspaces=true
}


\newlist{parts}{enumerate}{2}
\setlist[parts,1]{label=(\alph*),leftmargin=*,itemsep=4pt,topsep=3pt}
\setlist[parts,2]{label=(\roman*),leftmargin=*,itemsep=2pt,topsep=2pt}

\newcommand{\pts}[1]{\hfill   \textit{[#1]}}


\begin{document}

\begin{center}
  {\large\bfseries BANARAS HINDU UNIVERSITY}\\[2pt]
  {\normalsize B.Voc. Semester IV Examination, 2022-23}\\[6pt]
  \rule{0.8\linewidth}{0.5pt}\\[6pt]
  {\large\bfseries Paper: BVCA-402 --- Probability and Statistics}\\[4pt]
  {\normalsize Time: Three Hours \hfill Maximum Marks: 70}
\end{center}

\rule{1.0\linewidth}{0.4pt}

\smallskip



\noindent   \textit{Instructions: Attempt any two questions from Section~A, any six questions from Section~B, and all questions from Section~C. The marks for each question are given against the respective questions.}

\bigskip

\bigskip


\noindent {\large\bfseries SECTION --- A}
\rule{\linewidth}{0.4pt}\smallskip

\medskip


\noindent   \textbf{Question 1.}

\begin{parts}
  \item State and prove Markov's inequality.\pts{3}
  \item Derive Chebyshev's inequality using Markov's inequality.\pts{3}
  \item Let $f(x)$ be the uniform distribution on $0 \le x \le 5$ and $0$ everywhere else. Give a bound using Chebyshev's inequality for $P(1 \le X \le 4)$ and also calculate the actual probability. How do they compare?\pts{6}
\end{parts}

\medskip


\noindent   \textbf{Question 2.}

\begin{parts}
  \item Give the axiomatic definition of probability.\pts{2}
  \item Define convergence in distribution.\pts{2}
  \item Consider the joint probability mass function of two discrete random variables $X$ and $Y$ as $f(x, y) = \frac{xy^2}{13}$ with support $\{(x, y) = (1,1), (1, 2), (2, 2)\}$. Construct a joint probability distribution table of $X$ and $Y$. Are $X$ and $Y$ independent? Find $E(X)$ and $E(Y)$.\pts{8}
\end{parts}

\medskip


\noindent   \textbf{Question 3.}

\begin{parts}
  \item Define simple and composite hypothesis with examples.\pts{2}
  \item What are Type-I and Type-II errors in hypothesis testing?\pts{2}
  \item Urban storm water can be contaminated by many sources, including discarded batteries. When ruptured, these batteries release metals of environmental significance. A sample of $49$ Panasonic AAA batteries gave a sample mean zinc mass of $2.06 \text{ g}$ and a sample standard deviation of $0.140 \text{ g}$. Does this data provide compelling evidence for concluding that the population mean zinc mass exceeds $2.0 \text{ g}$? Construct a $95\%$ confidence interval for the population mean zinc mass. Choose the appropriate critical value(s) ($z_{0.05} = 1.64, t_{0.05, 48} = 1.68, t_{0.025, 48} = 2.01, z_{0.025} = 1.96$).\pts{8}
\end{parts}

\medskip


\noindent   \textbf{Question 4.} Let $X$ be a random variable with PDF given by:
\[ f_X(x) = 
\begin{cases} x^2(2x+c) &  \text{if } 0 < x \le 1 \\ 0 &  \text{otherwise} \end{cases} \]
Find the value of $c$. Also compute the mean, variance and CDF of the random variable $X$. If $Y = \frac{2}{x} + 3$ is a function of the random variable $X$, then find $E(Y)$ and $\operatorname{Var}(Y)$.\pts{12}

\bigskip


\noindent {\large\bfseries SECTION --- B}
\rule{\linewidth}{0.4pt}\smallskip

\medskip


\noindent   \textbf{Question 5.}

\begin{parts}
  \item If the random variables $X$ and $Y$ independently follow exponential distribution with parameters $ \theta$ and $\lambda$, respectively, then find the expectation and variance of $Z = X + Y$.\pts{4}
  \item What do you mean by standardization of a variable?\pts{2}
\end{parts}

\medskip


\noindent   \textbf{Question 6.}

\begin{parts}
  \item Show that the Pearson's correlation coefficient between the variables $X$ and $Y$ is independent of change of scale and origin.\pts{4}
  \item Assume neither $A$ nor $B$ is an impossible event. If $A$ and $B$ are mutually exclusive, will they be independent?\pts{2}
\end{parts}

\medskip


\noindent   \textbf{Question 7.}

\begin{parts}
  \item Show that the Poisson distribution is a limiting case of the Binomial distribution.\pts{4}
  \item What is the p-value in hypothesis testing?\pts{2}
\end{parts}

\medskip


\noindent   \textbf{Question 8.} A person uses his car $30\%$ of the time, walks $30\%$ of the time and rides the bus $40\%$ of the time as he goes to work. He is late $10\%$ of the time when he walks; he is late $3\%$ of the time when he drives; and he is late $7\%$ of the time when he takes the bus. What is the probability he took the bus if he was late? What is the probability he walked if he is on time? What is the probability that he'll be on time?\pts{6}

\medskip


\noindent   \textbf{Question 9.} Vehicles pass through a junction on a busy road at an average rate of $240$ per hour. Find the probability that none passes in a given minute. What is the expected number passing in three minutes? Find the probability that this expected number actually passes through in a given three-minute period. (Assume $e \approx 2.72$).\pts{6}

\medskip


\noindent   \textbf{Question 10.}

\begin{parts}
  \item State and prove the law of total probability.\pts{4}
  \item Give an example where two events are mutually exclusive but not exhaustive. Also give an example where two events are not mutually exclusive but exhaustive.\pts{2}
\end{parts}

\medskip


\noindent   \textbf{Question 11.} Define Normal distribution. Also discuss its important properties.\pts{6}

\medskip


\noindent   \textbf{Question 12.}

\begin{parts}
  \item Suppose the average CGPA of students is normally distributed with population mean $8.7$ and standard deviation $0.2$. At $10\%$ level of significance, what is the minimum sample size needed to obtain a margin of error of $0.013$?\pts{4}
  \item What is a cumulative distribution function? What are the criteria that a valid CDF needs to satisfy?\pts{2}
\end{parts}

\medskip


\noindent   \textbf{Question 13.} The results of a national survey showed that on average, adults sleep $6.9$ hours per night. Suppose that the standard deviation is $1.2$ hours. Use Chebyshev's theorem to calculate the minimum percentage of individuals who sleep between $4.5$ and $9.3$ hours. Use Chebyshev's theorem to calculate the minimum percentage of individuals who sleep between $3.9$ and $9.9$ hours. Assume that the number of hours of sleep follows a normal distribution. Use the empirical rule to calculate the percentage of individuals who sleep between $4.5$ and $9.3$ hours per day. How does this result compare to the value that you obtained using Chebyshev's theorem in the first part of this question?\pts{6}

\bigskip


\noindent {\large\bfseries SECTION --- C}
\rule{\linewidth}{0.4pt}\smallskip

\medskip


\noindent   \textbf{Question 14.} Answer all the questions in a sentence each:\pts{10x1=10}

\begin{parts}
  \item What is the sampling distribution of sample variance?
  \item What is the range of Pearson's correlation coefficient?
  \item When will the central limit theorem fail, even if the sample size is large?
  \item What is the relation between power of the test and Type-II error?
  \item Construct a sample space for the event of tossing a coin till the first head appears.
  \item Let $P(B) > 0$ so that $P(A|B)$ is defined. If $A \subseteq B$, what is $P(A|B)$?
  \item Define probability mass function with a suitable example.
  \item How can we get the binomial distribution from Bernoulli random variables?
  \item If mean = median = mode, what can we say about that distribution?
  \item What is the formula for sample covariance?
\end{parts}

\vfill
\rule{\linewidth}{0.4pt}

\begin{center}\small
\href{https://bhu-exam.vercel.app/}{Click here to visit our website}\\[4pt]\href{https://me-aryan.vercel.app/}{Click here to visit our contributor}\\[4pt]\href{https://quashan-me.vercel.app/}{Click here to visit our contributor}
\end{center}
\end{document}
